% Corrigé intégré automatiquement pour la banque Exercices.
% Source : A_Integrer/DDS_02/042_Chs_Leq/042_Chs_Leq.tex
\begin{corrigebox}[Corrigé]
\CorrigeQuestion{1}
On paramètre les mouvements relatifs, puis on applique la composition des vitesses et, si nécessaire, la dérivation vectorielle dans le repère adapté. Pour un train épicycloïdal, on écrit la relation de Willis avec les nombres de dents avant d'imposer les éléments bloqués ou entraînés.
\CorrigeQuestion{2}
On identifie les classes d'équivalence cinématique, les surfaces de contact et les mouvements relatifs autorisés. Chaque contact est remplacé par la liaison normalisée correspondante, puis les axes et points caractéristiques sont reportés sur le schéma cinématique minimal.
\end{corrigebox}
